% Three shapes a closed walk can trace: a tree, a cycle, and two cycles sharing a vertex.
% Requires: \shapeunit (unit length), \shapesep (gap between shapes), colours ink and accent.
\begin{tikzpicture}[x=\shapeunit, y=\shapeunit,
    v/.style={circle, fill=ink, inner sep=1.6pt}, e/.style={ink, semithick}]
  % tree: v = 5, e = 4
  \begin{scope}[yshift={-1*\shapeunit}]
    \node[v] (a) at (0,2) {}; \node[v] (b) at (-1,1) {}; \node[v] (c) at (1,1) {};
    \node[v] (d) at (-1.5,0) {}; \node[v] (f) at (-0.5,0) {};
    \draw[e] (a)--(b) (a)--(c) (b)--(d) (b)--(f);
    \node[font=\footnotesize, align=center, anchor=north] at (0,-0.45) {tree\\ $v-e=1$\\ $\asymp n$ copies};
  \end{scope}
  % cycle: v = 5, e = 5
  \begin{scope}[xshift=\shapesep]
    \foreach \i in {0,...,4} \node[v] (p\i) at ({90+72*\i}:1.05) {};
    \draw[e] (p0)--(p1)--(p2)--(p3)--(p4)--(p0);
    \node[font=\footnotesize, align=center, anchor=north] at (0,-1.45) {one cycle\\ $v-e=0$\\ $\asymp 1$ copies};
  \end{scope}
  % two triangles sharing a vertex: v = 5, e = 6
  \begin{scope}[xshift={2*\shapesep}]
    \node[v, fill=accent] (o) at (0,0) {};
    \node[v] (l1) at (-1.3,0.75) {}; \node[v] (l2) at (-1.3,-0.75) {};
    \node[v] (r1) at (1.3,0.75) {};  \node[v] (r2) at (1.3,-0.75) {};
    \draw[e] (o)--(l1)--(l2)--(o) (o)--(r1)--(r2)--(o);
    \node[font=\footnotesize, align=center, anchor=north] at (0,-1.45) {two cycles\\ $v-e=-1$\\ $\asymp 1/n$ copies};
  \end{scope}
\end{tikzpicture}
